\section*{الفصل \ref{ch:b3:homogeneous-catalysis} --- الحفز المتجانس في الصناعة}

\begin{solution}{exo:b3:homogeneous-catalysis:1}
$\ce{RhCl(PPh3)3}$: 16، Rh(I) $d^8$. $\ce{RhCl(PPh3)2}$: 14، Rh(I). $\ce{RhClH2(PPh3)2}$: 16، Rh(III) $d^6$.
معقد الألكين: 18، Rh(III). هيدريد الألكيل: 16، Rh(III). والحذف الاختزالي
يعيد إلى 14، Rh(I).
\end{solution}

\begin{solution}{exo:b3:homogeneous-catalysis:2}
البوتانال، $\ce{CH3CH2CH2CHO}$ (الخطي)، و2-ميثيل بروبانال، $\ce{(CH3)2CHCHO}$ (المتفرع).
\end{solution}

\begin{solution}{exo:b3:homogeneous-catalysis:3}
سيكلوبنت-3-إين-1،1-ثنائي كربوكسيلات ثنائي الإيثيل، حلقة خماسية، مع فقد
الإيثين: تغادر مجموعتا $\ce{CH2}$ الطرفيتان في صورة $\ce{C2H4}$ وتتصل ذرتا الكربون CH الداخليتان.
\end{solution}

\begin{solution}{exo:b3:homogeneous-catalysis:4}
(\textit{E})-3-فينيل بروب-2-إينوات الميثيل (سينامات الميثيل): يُضاف الأريل إلى
الكربون الطرفي، ويعطي حذف الهيدريد $\beta$، بإزالة من الجهة نفسها، انطلاقًا من المطابق الأكثر استقرارًا
الألكين \textit{E}.
\end{solution}

\begin{solution}{exo:b3:homogeneous-catalysis:5}
الإضافة التأكسدية بطيئة و$\ce{[Rh(CO)2I2]-}$ هو حالة الراحة، ومنه
$v = k[\mathrm{Rh}]_{\mathrm{tot}}[\ce{CH3I}]$. والميثانول لا يفعل إلا أن يجدد يوديد الميثيل بتفاعل سريع مع HI؛
واليوديد الكلي المشحون يحدد $[\ce{CH3I}]$، فلا تتوقف السرعة على تركيز
الميثانول حتى لا يبقى منه ما يكفي لإعادة تحويل HI إلى
يوديد الميثيل.
\end{solution}

\begin{solution}{exo:b3:homogeneous-catalysis:6}
حمض فينيل البورونيك مع 4-بروموتولوين، أو حمض 4-ميثيل فينيل البورونيك مع
البروموبنزين. وكلاهما يعمل؛ والأول يستعمل أرخص حمض بورونيك وبروميد أريل
شائعًا ومستقرًا.
\end{solution}

\begin{solution}{exo:b3:homogeneous-catalysis:7}
ذو التناظر $C_2$: أيزوتاكتي (يختار الموقعان كلاهما الوجه نفسه). ذو التناظر $C_s$:
سنديوتاكتي (موقعان هما صورتان مرآويتان يتناوبان). $\ce{Cp2ZrCl2}$ غير الجسري: أتاكتي (مواقع
لاكيرالية).
\end{solution}

\begin{solution}{exo:b3:homogeneous-catalysis:8}
$\mathrm{TON} = 0.98/0.0010 = 980$؛ ومتوسط $\mathrm{TOF} = 980/\qty{2.0}{h} = \qty{490}{h^{-1}}$.
\end{solution}

\begin{solution}{exo:b3:homogeneous-catalysis:9}
الوحدة المتكررة هي $\ce{-[CH=CH-C5H8]-}$ مع مجموعتي CH على الكربونين 1 و3 من
حلقة سيكلوبنتان (بولي(1،3-سيكلوبنتيلين فينيلين)). وتحرر
إجهاد الحلقة الثنائية يجعل $\Delta_rH^\circ$ سالبة بشدة، وهذا يرجح على
فقد الإنتروبيا الانتقالية؛ ولا يتكوّن أي غاز، فالقوة الدافعة
أنتالبية.
\end{solution}

\begin{solution}{exo:b3:homogeneous-catalysis:10}
$\bar n_n = 1/(1 - p) = 2000$؛ $M_w/M_n = 1 + p = 1.9995 \approx 2.0$. ولحفاز تسيغلر--ناتا
أنواع متعددة من المواقع، لكل منها $p$ الخاص به؛ ومجموع عدة توزيعات شولتس--فلوري
بمتوسطات مختلفة أعرض (التشتت غالبًا من 4 إلى 8).
\end{solution}

\begin{solution}{exo:b3:homogeneous-catalysis:11}
عند الضغط المنخفض $K_1K_2[\ce{H2}] \ll [\mathrm L] + K_1$: $v \approx kK_1K_2[\mathrm{Rh}]_{\mathrm{tot}}[\text{الألكين}][\ce{H2}]/([\mathrm L] + K_1)$،
من الرتبة الأولى بالنسبة إلى $\ce{H2}$. وعند الضغط المرتفع يغلب الحد $K_1K_2[\ce{H2}]$: $v \to k[\mathrm{Rh}]_{\mathrm{tot}}[\text{الألكين}]$،
مستقلة عن $\ce{H2}$. ولا يظهر $[\mathrm L]$ إلا في المقام: فإضافة الفوسفين تخفض $v$
بدفع الروديوم عائدًا إلى طليعة الحفاز المشبعة.
\end{solution}

\begin{solution}{exo:b3:homogeneous-catalysis:12}
الإدخال الذي يضع الروديوم على الكربون الداخلي يبني ألكيلًا ثانويًا
بجوار فلز مزدحم؛ والفوسفينات الضخمة (وفائض منها، يُبقي اثنين
منها على الفلز) تجعل تلك الحالة الانتقالية أسوأ من الحالة التي تعطي الألكيل
الأولي، الذي يقود إلى الألدهيد الخطي. والبوتانال الخطي هو الذي يحوَّل
بالتكاثف الألدولي والهدرجة إلى الكحول المتفرع $C_8$ المستعمل في
الملدِّنات؛ أما 2-ميثيل بروبانال فناتج جانبي أدنى قيمة.
\end{solution}

\begin{solution}{pb:b3:homogeneous-catalysis:1}
\textbf{1.} Ir(I)، $d^8$: $9 + 4 + 2 + 1 = 16$ إلكترونًا (العدّ المتعادل، مع الشحنة).

\textbf{2.} $\ce{[CH3Ir(CO)2I3]-}$، 18 إلكترونًا، Ir(III).

\textbf{3.} $\ce{[CH3Ir(CO)3I2]}$، 18 إلكترونًا، Ir(III).

\textbf{4.} $\ce{[CH3COIr(CO)2I2]}$، 16 إلكترونًا، Ir(III).

\textbf{5.} $\ce{[CH3COIr(CO)2I3]-}$ (18) يحذف $\ce{CH3COI}$، فيتجدد $\ce{[Ir(CO)2I2]-}$.

\textbf{6.} $\ce{CH3COI + H2O -> CH3COOH + HI}$; $\ce{CH3OH + HI -> CH3I + H2O}$.

\textbf{7.} $\ce{CH3OH + CO -> CH3COOH}$.

\textbf{8.} مع الإيريديوم يكون الإدخال الهجري بطيئًا؛ أما الإضافة التأكسدية ليوديد
الميثيل، المحددة للسرعة مع الروديوم، فأسرع بكثير على الإيريديوم الأغنى
بالإلكترونات.

\textbf{9.} تنزع يوديدًا من $\ce{[CH3Ir(CO)2I3]-}$، فتعطي الثلاثي الكربونيل المتعادل، الذي
يكون فيه التبرع العكسي للفلز أقل وتهاجر فيه مجموعة الميثيل أسرع إلى CO أكثر
محبة للإلكترونات.

\textbf{10.} اعتماد عكسي على تركيز اليوديد الحر: فاليوديد يدفع
التوازن المسبق عائدًا نحو المعقد الأنيوني البطيء.

\textbf{11.} لا يدخل الميثانول إلا بعد الخطوة المحددة للسرعة، إذ يتحول إلى
يوديد الميثيل بتفاعل سريع مع HI.

\textbf{12.} $M(\ce{CH3COOH}) = \qty{60.1}{g/mol}$: $n = \qty{5.0e11}{g}/\qty{60.1}{g/mol} = \qty{8.3e9}{mol}$ في السنة.

\textbf{13.} $\qty{8.32e9}{mol}/\qty{8000}{h} = \qty{1.04e6}{mol/h}$.

\textbf{14.} $\qty{2.0e5}{L} \times \qty{5.0e-3}{mol/L} = \qty{1.0e3}{mol}$ من الإيريديوم، و$\qty{1.0e3}{mol} \times \qty{192.2}{g/mol} \approx \qty{190}{kg}$.

\textbf{15.} $\mathrm{TOF} = \qty{1.04e6}{mol/h}/\qty{1.0e3}{mol} = \qty{1.0e3}{h^{-1}}$، أي نحو \qty{0.29}{s^{-1}}.

\textbf{16.} $\mathrm{TON} = \qty{8.3e9}{mol}/\qty{1.0e3}{mol} = \num{8.3e6}$ في السنة.

\textbf{17.} $\qty{5.0e8}{kg}/(\qty{8000}{h} \times \qty{200}{m^3}) \approx \qty{310}{kg.m^{-3}.h^{-1}}$.

\textbf{18.} $\ce{CO + H2O -> CO2 + H2}$.

\textbf{19.} $1/0.94 = \qty{1.064}{mol}$ من CO لكل مول من حمض الأسيتيك.

\textbf{20.} $1.064 - 1 = \qty{0.064}{mol}$ من \ce{CO2} لكل مول من حمض الأسيتيك.

\textbf{21.} الطن الواحد يساوي $\qty{1.0e6}{g}/\qty{60.1}{g/mol} = \qty{1.66e4}{mol}$؛ $\qty{1.66e4}{mol} \times 0.0638 \times \qty{44.0}{g/mol} \approx \qty{47}{kg}$ من \ce{CO2}.

\textbf{22.} الماء متفاعل في تفاعل الإزاحة: فتقليله يبطّئها ويوفّر طاقة
تجفيف الحمض. لكن الماء يظل لازمًا لحلمأة يوديد الأسيتيل ولإبقاء
الحفاز نشطًا ومذابًا؛ ومع الروديوم، يؤدي نقص الماء إلى ترسيب يوديد
الروديوم.

\textbf{23.} الهيدروجين: يتراكم في الطور الغازي، ويخفض الضغط الجزئي للغاز CO
(فيجب تصريف الغاز، مع فقد CO)، ويمكن أن يهدرج الوسائط فيعطي
نواتج جانبية مثل حمض البروبيونيك.

\textbf{24.} يدور حفاز الإيريديوم نحو $\qty{1.0e3}{h^{-1}}$: فكل ذرة إيريديوم تصنع
ألف جزيء من حمض الأسيتيك في الساعة، أي نحو ثمانية ملايين في السنة.
\end{solution}
